% =====================================================================
% chapters/06-conclusion.tex
% PT: Capítulo 6 -- Conclusão. Exemplo de: resumo das contribuições,
%     respostas às questões de investigação e trabalho futuro (enumerate).
% EN: Chapter 6 -- Conclusion. Shows: summary of the contributions,
%     answers to the research questions and future work (enumerate).
% =====================================================================
\chapter{Conclusion}
\label{chap:conclusion}

% PT: Parágrafo de abertura do capítulo.
% EN: Chapter opening paragraph.
This chapter closes the document. In a real dissertation, the conclusion
summarises what was done, answers the research questions, states the
contributions and suggests future work. It should not introduce new results.

% ---------------------------------------------------------------------
% PT: Exemplo em inglês
% EN: English example
% ---------------------------------------------------------------------
\section{English example}
\label{sec:conclusion:en}

\subsection{Summary and Contributions}
\label{sec:conclusion:summary}

Answer each research question from \cref{sec:intro:questions} in turn, citing
the section, table or figure that supports the answer, for example
\cref{tab:results}. Then list the main contributions of the work in one or
two short paragraphs.

\subsection{Future Work}
\label{sec:conclusion:future}

\begin{enumerate}
  \item Evaluate the proposed method on a larger and more varied dataset.
  \item Compare it with deeper architectures, such as residual networks
        \cite{He2016}.
  \item Study the cost of deploying the method in a real setting.
\end{enumerate}

% ---------------------------------------------------------------------
% PT: Exemplo em português
% EN: Portuguese example
% ---------------------------------------------------------------------
\section{Exemplo em português}
\label{sec:conclusion:pt}

\begin{otherlanguage}{portuguese}

A conclusão responde às questões de investigação da
Secção~\ref{sec:intro:questions}, resume os contributos do trabalho e indica
direções de trabalho futuro. Não deve apresentar resultados novos nem
repetir a discussão do Capítulo~\ref{chap:discussion}.

\end{otherlanguage}
